\documentclass[10pt,a4paper]{amsart}
\usepackage[margin=19mm]{geometry}
\usepackage{amsmath,amssymb,mathtools}
\usepackage[scr=rsfs,cal=euler]{mathalfa}
\usepackage{xfrac}
\usepackage[unicode,hidelinks,bookmarks=false]{hyperref}
\newcommand{\ie}{i.e.\ }
\newcommand{\cf}{cf.\ }
\newcommand{\resp}{respectively\ }
\newcommand{\Subsection}{\subsection}
\providecommand{\nobreakdash}{\nobreak\hbox{-}}
\setlength{\emergencystretch}{2em}
\multlinegap0pt
\usepackage[shortlabels]{enumitem}
\usepackage{algorithm}\usepackage{algpseudocode}
\usepackage{amssymb}
\usepackage{xcolor}
\usepackage{tikz}
\newtheorem{assumption}{Assumption}
\newtheorem{lemma}{Lemma}
\newcommand{\R}{\mathbb{R}}
\newcommand{\norm}[1]{\lVert #1\rVert}
\newcommand{\tc}[1]{#1}
\title{Dirac--Frenkel dynamics with inertia for nonlinearly parametrized solutions of evolution problems}
\author{Matteo Raviola, Benjamin Peherstorfer}
\date{Source: arXiv:2606.24769v3 (CC BY 4.0)}
\begin{document}
\maketitle

\noindent\emph{Source and licence.} Dirac--Frenkel dynamics with inertia for nonlinearly parametrized solutions of evolution problems. Version arXiv:2606.24769v3 (CC BY 4.0), distributed by arXiv under the Creative Commons Attribution 4.0 International licence, http://creativecommons.org/licenses/by/4.0/, as recorded in the arXiv licence metadata for this exact version and shown on the abstract page https://arxiv.org/abs/2606.24769v3. Full licence notice: \texttt{licenses/arxiv-2606.24769-CC-BY-4.0.txt}.\par\smallskip

\noindent\textit{Editorial scope.} Counted unit: the article's lemma lem:local\_error\_new (source0.tex lines 914--938). The article's complete original proof of it is delivered with it (source0.tex lines 939--1117, one \texttt{proof} environment of 179 lines). No mathematics is written, repaired or re-derived: the statement and the proof are the article's own lines, in the article's own notation, and no retained line is substituted or re-flowed.  Retained uncounted as the source-local context the proof needs: lines 173--178; lines 847--875; lines 879--888.  Nothing was omitted from any retained span, so the package carries no omission record.  No retained line contains a citation, so the wrapper carries no bibliography and no bibliography entry.  No retained line contains a figure environment, an image-inclusion command or drawing code, so no image is delivered and the package declares no asset dependency.  Editorial formatting only, in addition: an AMS \texttt{amsart} preamble that restates the article's own declarations and macros used by the retained lines, namely the environments \texttt{assumption}, \texttt{lemma} on separate counters, together with the standard packages those lines need.  The retained environments print in this document as Assumption 1, Lemma 1 in order of appearance, whereas the article prints the same items with the numbering of its own full document.  The complete native source of the article as distributed in the arXiv e-print for this exact version is bundled unchanged as \texttt{sources/203547/source0.tex}.
        \begin{equation}
          \hat u=\Phi(\theta),
          \qquad J(\theta)w:=D\Phi(\theta)[w]\quad(w\in \Theta),
          \qquad f(\theta):=F(\hat u)\in H.
          \label{eq:Jf_def}
        \end{equation}

\begin{assumption}[Flow stability and regularity]
  \label{ass:flow_stability_new}
  Fix $T>0$.
  Assume that there exist constants $\ell\in\R$, $C_\Phi\ge0$, $c_\Phi\ge0$, and $C_a\ge0$ such that the following hold for all relevant states.
  \begin{enumerate}[label=\textnormal{(\roman*)},leftmargin=2.5em]
    \item The exact flow satisfies
          \begin{equation}
            \norm{\varphi_t(u)-\varphi_t(\widetilde u)}_H
            \le e^{\ell t}\norm{u-\widetilde u}_H,
            \qquad 0\le t\le T.
            \label{eq:flow_stability_new}
          \end{equation}
    \item The second derivative of $\Phi$ is bounded, relative to the frozen Jacobian $J(\theta)w:=D\Phi(\theta)[w]$ at the base point, as follows for all relevant $\theta,\zeta,\xi$:
          \begin{equation}
            \norm{D^2\Phi(\theta+s\zeta)[\xi,\xi]}_H
            \le
            C_\Phi\norm{J(\theta)\xi}_H^2
            +
            c_\Phi\norm{\xi}_\Theta^2,
            \qquad 0\le s\le 1.
            \label{eq:D2Phi_bound_new}
          \end{equation}
    \item For every relevant initial state $w$ and the corresponding exact trajectory $y(s)=\varphi_s(w)$, one has
          \[
            \norm{\ddot y(s)}_H\le C_a
          \]
          for all times $s$ for which the trajectory is used below.
  \end{enumerate}
\end{assumption}

\begin{equation}
  \Delta_k^2
  :=
  \norm{J_kv_{k+1}-f_k}_H^2
  +
  \varepsilon^2\norm{v_{k+1}}_\Theta^2
  +
  \frac{\tau^2}{h_k}\norm{v_{k+1}-v_k}_\Theta^2\,.
  \label{eq:Delta_min_error}
\end{equation}

\begin{lemma}[One-step local error]
  \label{lem:local_error_new}
  Under Assumption~\ref{ass:flow_stability_new}, for every time step $k$,
  the local error of one Euler step
  \begin{equation}\label{eq:one_step_uplus}
    \hat{u}_+ := \Phi(\theta_k + h_kv_{k + 1})\,.
  \end{equation}
  satisfies
  \begin{equation}
    \norm{\hat u_+-\varphi_{h_k}(\hat u_k)}_H
    \le
    h_k\,\Delta_k
    +
    C_\Phi h_k^2\norm{f_k}_H^2
    +
    \tc{
    h_k^2
    \left(
    C_\Phi+\frac{c_\Phi}{2\varepsilon^2}
    \right)\Delta_k^2}
    +
    \frac{C_a}{2}h_k^2 .
    \label{eq:local_error_bound_new}
  \end{equation}
\end{lemma}

\begin{proof}
  Let $ y_k(s):=\varphi_s(\hat u_k)$ for $0\le s\le h_k$.
  Then $y_k(0)=\hat u_k$ and, since $y_k$ solves the exact evolution equation, $\dot y_k(0)=F(\hat u_k)=f_k$.
  Taylor's formula with integral remainder gives
  \begin{equation}
    \varphi_{h_k}(\hat u_k)
    =
    y_k(h_k)
    =
    y_k(0)+h_k\dot y_k(0)
    +
    \int_0^{h_k}(h_k-s)\ddot y_k(s)\,ds .
    \label{eq:flow_taylor_integral_new}
  \end{equation}
  Hence
  \begin{equation}
    \varphi_{h_k}(\hat u_k)
    =
    \hat u_k+h_k f_k+r_k^{\mathrm{flow}},
    \label{eq:flow_taylor_new}
  \end{equation}
  where
  \begin{equation}
    r_k^{\mathrm{flow}}
    :=
    \int_0^{h_k}(h_k-s)\ddot y_k(s)\,ds .
    \label{eq:flow_remainder_def_new}
  \end{equation}
  By Assumption~\ref{ass:flow_stability_new}\textnormal{(iii)},
  \begin{align}
    \norm{r_k^{\mathrm{flow}}}_H
     & \le
    \int_0^{h_k}(h_k-s)\norm{\ddot y_k(s)}_H\,ds
    \notag \\
     & \le
    C_a\int_0^{h_k}(h_k-s)\,ds
    =
    \frac{C_a}{2}h_k^2 .
    \label{eq:flow_remainder_bound_new}
  \end{align}
  Next, Taylor's formula for $\Phi$ applied to the curve $s \mapsto \theta_k + sh_k v_{k + 1}$, which parametrizes the straight line segment in $\Theta$ from $\theta_k$ to $\theta_{k + 1} = \theta_k + h_k v_{k + 1}$,  gives
  \begin{align}
    \hat u_+
     & =
    \Phi(\theta_k+h_kv_{k+1})
    \notag \\
     & =
    \Phi(\theta_k)
    +
    h_kD\Phi(\theta_k)[v_{k+1}]
    +
    h_k^2
    \int_0^1
    (1-s)
    D^2\Phi(\theta_k+s h_kv_{k+1})
    [v_{k+1},v_{k+1}]
    \,ds .
    \label{eq:Phi_taylor_integral_new}
  \end{align}
  Since $\hat u_k=\Phi(\theta_k)$ and
  $J_kv_{k+1}=D\Phi(\theta_k)[v_{k+1}]$ (see \eqref{eq:Jf_def}), this becomes
  \begin{equation}
    \hat u_+
    =
    \hat u_k+h_kJ_kv_{k+1}+r_k^\Phi,
    \label{eq:Phi_taylor_new}
  \end{equation}
  where
  \begin{equation}
    r_k^\Phi
    :=
    h_k^2
    \int_0^1
    (1-s)
    D^2\Phi(\theta_k+s h_kv_{k+1})
    [v_{k+1},v_{k+1}]
    \,ds .
    \label{eq:Phi_remainder_def_new}
  \end{equation}
  Applying Assumption~\ref{ass:flow_stability_new}\textnormal{(ii)}
  with $
    \theta=\theta_k,
    \zeta=h_kv_{k+1},
    \xi=v_{k+1}$
  yields
  \begin{align}
    \norm{r_k^\Phi}_H
     & \le
    h_k^2
    \int_0^1
    (1-s)
    \norm{
    D^2\Phi(\theta_k+s h_kv_{k+1})
    [v_{k+1},v_{k+1}]
    }_H
    \,ds
    \notag \\
     & \le
    h_k^2
    \int_0^1
    (1-s)
    \left(
    C_\Phi\norm{J_kv_{k+1}}_H^2
    +
    c_\Phi\norm{v_{k+1}}_\Theta^2
    \right)
    \,ds
    \notag \\
     & =
    \frac{h_k^2}{2}
    \left(
    C_\Phi\norm{J_kv_{k+1}}_H^2
    +
    c_\Phi\norm{v_{k+1}}_\Theta^2
    \right).
    \label{eq:Phi_remainder_bound_new}
  \end{align}
  Let us now bound the state velocity $J_kv_{k+1}$.
    \tc{
      Since
      \begin{equation}
        J_kv_{k+1}=f_k+(J_kv_{k+1}-f_k),
      \end{equation}
      the definition of $\Delta_k$ gives
      \begin{equation} \label{eq:state_velocity_bound}
        \norm{J_kv_{k+1}}_H^2
        \le
        2\norm{f_k}_H^2
        +
        2\norm{J_kv_{k+1}-f_k}_H^2
        \le
        2\norm{f_k}_H^2
        +
        2\Delta_k^2.
      \end{equation}
      Moreover,
      \begin{equation} \label{eq:velocity_bound}
        \norm{v_{k+1}}_\Theta^2
        \le
        \frac{\Delta_k^2}{\varepsilon^2}.
      \end{equation}
    }

  We now subtract the exact-flow expansion
  \eqref{eq:flow_taylor_new} from the parametric expansion
  \eqref{eq:Phi_taylor_new} and obtain with the triangle inequality,
  \begin{align}
    \norm{\hat u_{k+1}-\varphi_{h_k}(\hat u_k)}_H
     & \le
    h_k\norm{J_kv_{k+1}-f_k}_H
    +
    \norm{r_k^\Phi}_H
    +
    \norm{r_k^{\mathrm{flow}}}_H .
    \label{eq:local_difference_bound_new}
  \end{align}
  Inserting the bounds
  \eqref{eq:flow_remainder_bound_new},
  \eqref{eq:Phi_remainder_bound_new},
  \tc{ \eqref{eq:state_velocity_bound},
      and \eqref{eq:velocity_bound}}, we obtain
  \begin{align}
    \norm{\hat u_{k+1}-\varphi_{h_k}(\hat u_k)}_H
     & \le
    h_k\norm{J_kv_{k+1}-f_k}_H
    +
    C_\Phi h_k^2\norm{f_k}_H^2
    +
    \tc{
    h_k^2
    \left(
    C_\Phi+\frac{c_\Phi}{2\varepsilon^2}
    \right)\Delta_k^2}
    +
    \frac{C_a}{2}h_k^2 .
    \label{eq:local_error_before_defect_new}
  \end{align}
  Finally, by using \eqref{eq:Delta_min_error}, we obtain the bound \eqref{eq:local_error_bound_new}.
\end{proof}

\end{document}
